\documentclass[11pt]{article}
\usepackage{amsmath, amssymb}

\author{Vinu Vikraman}
\date{}

\begin{document}


\begin{center}
    {\huge \textbf{Description of input.ini keywords}}
\end{center}

%\section{}

\section{\texttt{[imagecata]}}
\subsection{\texttt{imagefile}}
\begin{itemize}
\item Specifies the main input image file.
\item This is typically a FITS file containing the science data to be analyzed.
\item Example:
\begin{verbatim}
imagefile = frame-r-002125-3-0104.fits
\end{verbatim}
\end{itemize}

\subsection{\texttt{whtfile}}
\begin{itemize}
\item Specifies the weight or noise map associated with the input image.
\item The interpretation depends on the filename:
\begin{itemize}
\item If it contains rms'', it is treated as an RMS map. \item If it contains weight'', it is treated as a weight map.
\item Otherwise, it is assumed to be an RMS map by default.
\end{itemize}
\item Example:
\begin{verbatim}
whtfile = sigma.fits
\end{verbatim}
\end{itemize}

\subsection{\texttt{sex\_cata}}
\begin{itemize}
\item Specifies the catalogue generated by a source detection algorithm.
\item Contains detected objects along with their measured properties such as position, flux, and shape parameters.
\item Used as the primary detection catalogue for further processing.
\item Example:
\begin{verbatim}
sex\_cata = sex.cat
\end{verbatim}
\end{itemize}

\subsection{\texttt{obj\_cata}}
\begin{itemize}
\item Specifies an external catalogue of objects, typically obtained from an online database.
\item The expected format includes positional information such as right ascension (RA) and declination (DEC), possibly split into multiple columns.
\item This catalogue is used to define target objects for further analysis.
\item Example:
\begin{verbatim}
obj\_cata = input.cat
\end{verbatim}

In the case of an image which is not cutout \verb|galcut=False|, then accepted headers are
\begin{verbatim}
OBJ_ID RA DEC Z STAR 

OBJ_ID RA1 RA2 RA3 DEC1 DEC2 DEC3 STAR Z 
\end{verbatim}
If you are giving PSF file names in the config.ini or giving a file which contains the PSF files, then you can give only

\begin{verbatim}
OBJ_ID RA DEC Z 

OBJ_ID RA1 RA2 RA3 DEC1 DEC2 DEC3 Z
\end{verbatim}

If you don't need to convert the structural parameters to physical parameters, you can give the header

\begin{verbatim}
OBJ_ID RA DEC 

OBJ_ID RA1 RA2 RA3 DEC1 DEC2 DEC3 
\end{verbatim}


\end{itemize}

\subsection{\texttt{rootname}}
\begin{itemize}
\item Defines the base name for all output files generated by the pipeline.
\item Output files such as catalogues and images will use this prefix.
\item Example:
\begin{verbatim}
rootname = cl1358
\end{verbatim}
\end{itemize}


\subsection{\texttt{datadir}}
\begin{itemize}
\item Specifies the directory containing input image files.
\item If this keyword is omitted or commented out, the program defaults to using the current working directory.
\item Example:
\begin{verbatim}
datadir = /Users/vinu/github/pymorph_chatgpt/data_large
\end{verbatim}
\end{itemize}

\subsection{\texttt{outdir}}
\begin{itemize}
\item Specifies the directory where output data products will be stored.
\item If this keyword is omitted or commented out, output files are written to the current working directory.
\item Example:
\begin{verbatim}
outdir = /home/vinu/WorkAll/TNG/morphology/decomposition/
\end{verbatim}
\end{itemize}

\section{\texttt{[external]}}

\subsection{\texttt{GALFIT\_PATH}}
\begin{itemize}
\item Defines the full path to the GALFIT executable.
\item Used by the pipeline to invoke GALFIT for image fitting tasks.
\item Example:
\begin{verbatim}
/usr/local/bin/galfit
\end{verbatim}
\end{itemize}

\subsection{\texttt{SEX\_PATH}}
\begin{itemize}
\item Specifies the full path to the SExtractor executable.
\item Used for source detection and catalogue generation.
\item Example:
\begin{verbatim}
/opt/homebrew/bin/sex
\end{verbatim}
\end{itemize}

\section{\texttt{[psf]}}

\subsection{\texttt{stargal\_prob}}
\begin{itemize}
\item Threshold probability used for star-galaxy classification.
\item Objects with classification probability above this value are treated as stars.
\item Also used as a selection criterion in fitting modes such as \texttt{fit\_and\_fit}.
\item Example:
\begin{verbatim}
stargal_prob = 0.7
\end{verbatim}
\end{itemize}

\subsection{\texttt{star\_size}}
\begin{itemize}
\item Defines the size of the PSF image.
\item The PSF cutout size is computed as a multiple of the semi-major axis (SMA) obtained from SExtractor.
\item The final PSF size is given by:
[
\text{PSF size} = \texttt{star\_size} $\times$ \text{SMA}
]
\item Example:
\begin{verbatim}
star_size = 20
\end{verbatim}
\end{itemize}

\subsection{\texttt{psflist}}
\begin{itemize}
\item Specifies a list of PSF images to be used in the analysis.
\item Each PSF file name encodes positional information in the format:
[
\texttt{psf\_RA1RA2RA3+DEC1DEC2DEC3.fits}
]
\item The pipeline extracts Right Ascension (RA) and Declination (Dec) from the filename and updates the PSF header accordingly.
\item Multiple PSF files are separated by commas.
\item Alternatively, if specified as \verb|@filename|, the PSF list will be read from an external file.
\item Example:
\begin{verbatim}
psflist = psf_1447218+0838538.fits, psf_1447219+0828392.fits, ...

psflist = @filename.csv
\end{verbatim}
\end{itemize}

\section{\texttt{[mask]}}

\subsection{\texttt{manual\_mask}}
\begin{itemize}
\item Enables or disables manual masking.
\item If set to 1, masking is controlled manually; otherwise automatic masking is used.
\end{itemize}

\subsection{\texttt{mask\_reg}}
\begin{itemize}
\item Defines the radius of the masking region.
\item Mask radius is given by:
[
R\_{\text{mask}} = \texttt{mask\_reg} $\times$ (\text{semi-major axis of neighbour})
]
\end{itemize}

\subsection{\texttt{thresh\_area}}
\begin{itemize}
\item Sets the threshold ratio of neighbour area relative to the target object.
\item Neighbours with area less than this factor times the object area are masked.
\end{itemize}

\subsection{\texttt{threshold}}
\begin{itemize}
\item Defines the overlap condition for masking.
\item Masking occurs when scaled semi-major axes of object and neighbour overlap.
\end{itemize}

\subsection{\texttt{no\_mask}}
\begin{itemize}
\item If set to 1, masking is completely disabled.
\end{itemize}

\section{\texttt{[size]}}

\subsection{\texttt{size\_list}}
\begin{itemize}
\item Controls cutout size and resizing behavior.
\item Format:
[
[\texttt{resize}, \texttt{imsize}, \texttt{fracrad}, \texttt{square}, \texttt{fixsize}]
]
\end{itemize}
\texttt{resize = 0} and \texttt{imsize = 1} when using \texttt{galcut = True}. 
\texttt{resize = 1} and \texttt{imsize = 0} when using \texttt{galcut = False}. Then it will generate size of the image by multipling \text{fracrad} with semimajor axis of that object 

\subsection{\texttt{searchrad}}
\begin{itemize}
\item Search radius in arcseconds for identifying neighbouring objects.
\end{itemize}

\subsection{\texttt{nearest\_neighbour}}
\begin{itemize}
\item Boolean flag indicating whether to consider the nearest neighbour during processing. Don't change this value
\end{itemize}

\section{\texttt{[cosmology]}}

\subsection{\texttt{pixelscale}}
\begin{itemize}
\item Pixel scale in arcseconds per pixel.
\item Used to convert image measurements into physical units.
\end{itemize}

\section{\texttt{[nonparams]}}

\subsection{\texttt{back\_extraction\_radius}}
\begin{itemize}
\item Radius used for background estimation.
\end{itemize}

\subsection{\texttt{angle}}
\begin{itemize}
\item Rotation angle (in degrees), typically used for asymmetry calculations.
\end{itemize}

\section{\texttt{[modes]}}

\subsection{\texttt{galcut}}
\begin{itemize}
\item If True, assumes cutout images are already provided.
\end{itemize}

\subsection{\texttt{decompose}}
\begin{itemize}
\item Enables running of the decomposition (e.g., fitting using GALFIT).
\end{itemize}

\subsection{\texttt{detail}}
\begin{itemize}
\item Enables random perturbation of initial fitting parameters.
\end{itemize}

\subsection{\texttt{number\_random}}
\begin{itemize}
\item Number of random realizations used when \texttt{detail=True}.
\end{itemize}

\subsection{\texttt{casgm}}
\begin{itemize}
\item Enables computation of non-parametric morphology parameters (CASGM).
\end{itemize}

\subsection{\texttt{find\_and\_fit}}
\begin{itemize}
\item If True, finds and fits all objects within a magnitude range.
\end{itemize}

\section{\texttt{[findfit]}}

\subsection{\texttt{maglim}}
\begin{itemize}
\item Defines faint and bright magnitude limits for object selection.
\end{itemize}

\section{\texttt{[mag]}}

\subsection{\texttt{mag\_zero}}
\begin{itemize}
\item Photometric zero point for magnitude calculation.
\end{itemize}

\subsection{\texttt{photo\_filter}}
\begin{itemize}
\item Name of the photometric filter used (e.g., SDSS\_r).
\end{itemize}

\section{\texttt{[galfit]}}

\subsection{\texttt{components}}
\begin{itemize}
\item Specifies galaxy components to fit (e.g., bulge, disk, bar).
\end{itemize}

\subsection{\texttt{devauc}}
\begin{itemize}
\item If True, uses de Vaucouleurs profile ($n=4$) for bulge.
\end{itemize}

\subsection{\texttt{fitting}}
\begin{itemize}
\item Flags indicating which components or parameters are free during fitting.
The first for fitting bulge centre, second for disk center, next is for sky, next is for centre of the bar
\end{itemize}

\subsection{\texttt{bbox}, \texttt{barbox}}
\begin{itemize}
\item Initial boxiness parameters for bulge and bar components. If bbox=-99 it will not fit bulge boxiness. Also, barbox=-99 will not find bar boxiness
\end{itemize}

\section{\texttt{[version]}}

\subsection{\texttt{galfitv}}
\begin{itemize}
\item Version of GALFIT used.
\end{itemize}

\subsection{\texttt{pymorphv}}
\begin{itemize}
\item Version of the pipeline.
\end{itemize}

\section{\texttt{[diagnosis]}}

\subsection{\texttt{chi2nu\_limit}}
\begin{itemize}
\item Upper limit of reduced chi-square for a good fit.
\end{itemize}

\subsection{\texttt{goodness\_limit}}
\begin{itemize}
\item Sigma threshold for statistical goodness of fit.
\end{itemize}

\subsection{\texttt{center\_deviation\_limit}}
\begin{itemize}
\item Maximum allowed deviation between initial and fitted center.
\end{itemize}

\subsection{\texttt{do\_plot}}
\begin{itemize}
\item Enables plotting of results.
\end{itemize}

\section{\texttt{[params\_limit]}}

\subsection{\texttt{NMag}}
\begin{itemize}
\item Magnitude variation range relative to initial estimate.
\end{itemize}

\subsection{\texttt{NRadius}}
\begin{itemize}
\item Maximum radius scaling factor relative to initial estimate.
\end{itemize}

\subsection{\texttt{LN}, \texttt{UN}}
\begin{itemize}
\item Lower and upper limits of S'ersic index.
\end{itemize}

\subsection{\texttt{center\_constrain}}
\begin{itemize}
\item Maximum allowed positional shift during fitting.
\end{itemize}

\section{\texttt{[sextractor]}}
\begin{itemize}
\item Used to determine the initial configuration parameters. Certain parameters, such as GAIN, must be adjusted for each telescope
\end{itemize}

\newpage

\begin{center}
    {\huge \textbf{Non-parametric parameters}}
\end{center}

\setcounter{section}{0}
Non-parametric morphology measures are widely used to quantify the structure of galaxies without assuming specific models. The most commonly used parameters are Concentration ($C$), Asymmetry ($A$), Clumpiness ($S$), Gini coefficient ($G$), and $M_{20}$.

\section{Concentration ($C$)}
The concentration index measures how centrally concentrated the light is:
\begin{equation}
C = 5 \log \left( \frac{r_{80}}{r_{20}} \right)
\end{equation}
where $r_{20}$ and $r_{80}$ are the radii containing 20% and 80% of the total light, respectively.

\section{Asymmetry ($A$)}
Asymmetry quantifies how symmetric a galaxy is under $180^\circ$ rotation:
\begin{equation}
A = \frac{\sum | I(x,y) - I_{180}(x,y) |}{\sum | I(x,y) |} - A_{\text{bkg}}
\end{equation}
where $I_{180}$ is the image rotated by $180^\circ$, and $A_{\text{bkg}}$ accounts for background asymmetry.

\section{Clumpiness ($S$)}
Clumpiness measures the fraction of light in high-frequency structures:
\begin{equation}
S = \frac{\sum | I(x,y) - I_{S}(x,y) |}{\sum | I(x,y) |} - S_{\text{bkg}}
\end{equation}
where $I_S$ is the smoothed image.

\section{Gini Coefficient ($G$)}
The Gini coefficient measures the inequality in the distribution of pixel fluxes:
\begin{equation}
G = \frac{1}{\bar{X} n (n-1)} \sum_{i=1}^{n} (2i - n - 1) X_i
\end{equation}
where $X_i$ are the pixel fluxes sorted in increasing order, $\bar{X}$ is the mean flux, and $n$ is the number of pixels.

\section{$M_{20}$}
$M_{20}$ is the normalized second-order moment of the brightest 20% of the light:
\begin{equation}
M_{20} = \log_{10} \left( \frac{\sum_i M_i}{M_{\text{tot}}} \right)
\end{equation}
where the sum is over the brightest pixels containing 20% of the total flux, and
\begin{equation}
M_i = f_i \left[ (x_i - x_c)^2 + (y_i - y_c)^2 \right]
\end{equation}
\begin{equation}
M_{\text{tot}} = \sum_i M_i
\end{equation}
Here $f_i$ is the flux in pixel $i$, and $(x_c, y_c)$ is the galaxy center.

\section{Conclusion}
These parameters together provide a powerful, model-independent description of galaxy morphology and are widely used in galaxy evolution studies.


\end{document}

